%%%%% PROJECT-SPECIFIC MACROS %%%%%
% Define your project-specific terms here
% Import with: %%%%% PROJECT-SPECIFIC MACROS %%%%%
% Define your project-specific terms here
% Import with: %%%%% PROJECT-SPECIFIC MACROS %%%%%
% Define your project-specific terms here
% Import with: %%%%% PROJECT-SPECIFIC MACROS %%%%%
% Define your project-specific terms here
% Import with: \input{commands/macros}
%
% INSTRUCTIONS:
% 1. Use \textsc{} for method/dataset/model names (small caps)
% 2. Add \xspace at the end for automatic spacing
% 3. Group related terms together
% 4. Use consistent naming conventions

\usepackage{xspace}

%%%%%%%%%%%%%%%%%%%%%%%%%%%%%%%%%%%%%%%%%%%%%%%%%%%%%%%%%%%%%%%%%
% METHOD NAMES
%%%%%%%%%%%%%%%%%%%%%%%%%%%%%%%%%%%%%%%%%%%%%%%%%%%%%%%%%%%%%%%%%

% Your method name (replace with actual name)
% Example: \newcommand{\ours}{{\bf YourMethod}\xspace}
% Example: \newcommand{\methodname}{\textsc{MethodName}\xspace}

%%%%%%%%%%%%%%%%%%%%%%%%%%%%%%%%%%%%%%%%%%%%%%%%%%%%%%%%%%%%%%%%%
% DATASET NAMES
%%%%%%%%%%%%%%%%%%%%%%%%%%%%%%%%%%%%%%%%%%%%%%%%%%%%%%%%%%%%%%%%%

% Dataset names (use \textsc for consistency)
% Example: \newcommand{\datasetone}{\textsc{Dataset-One}\xspace}
% Example: \newcommand{\datasettwo}{\textsc{Dataset-Two}\xspace}

%%%%%%%%%%%%%%%%%%%%%%%%%%%%%%%%%%%%%%%%%%%%%%%%%%%%%%%%%%%%%%%%%
% MODEL NAMES
%%%%%%%%%%%%%%%%%%%%%%%%%%%%%%%%%%%%%%%%%%%%%%%%%%%%%%%%%%%%%%%%%

% LLM and model names
% Example: \newcommand{\gpt}{\textsc{GPT-4o-mini}\xspace}
% Example: \newcommand{\llama}{\textsc{Llama-3.1-70B-Instruct}\xspace}
% Example: \newcommand{\claude}{\textsc{Claude-3.5-Sonnet}\xspace}

%%%%%%%%%%%%%%%%%%%%%%%%%%%%%%%%%%%%%%%%%%%%%%%%%%%%%%%%%%%%%%%%%
% BASELINE NAMES
%%%%%%%%%%%%%%%%%%%%%%%%%%%%%%%%%%%%%%%%%%%%%%%%%%%%%%%%%%%%%%%%%

% Baseline method names
% Example: \newcommand{\baseline}{\textsc{BaselineMethod}\xspace}
% Example: \newcommand{\fewshot}{\textsc{Few-shot}\xspace}

%%%%%%%%%%%%%%%%%%%%%%%%%%%%%%%%%%%%%%%%%%%%%%%%%%%%%%%%%%%%%%%%%
% ABBREVIATIONS
%%%%%%%%%%%%%%%%%%%%%%%%%%%%%%%%%%%%%%%%%%%%%%%%%%%%%%%%%%%%%%%%%

% Common abbreviations used in your paper
% Example: \newcommand{\ie}{i.e.,\xspace}
% Example: \newcommand{\eg}{e.g.,\xspace}
% Example: \newcommand{\cf}{cf.\xspace}
% Example: \newcommand{\wrt}{w.r.t.\xspace}

%%%%%%%%%%%%%%%%%%%%%%%%%%%%%%%%%%%%%%%%%%%%%%%%%%%%%%%%%%%%%%%%%
% DOMAIN-SPECIFIC TERMS
%%%%%%%%%%%%%%%%%%%%%%%%%%%%%%%%%%%%%%%%%%%%%%%%%%%%%%%%%%%%%%%%%

% Terms specific to your research domain
% Keep consistent formatting throughout the paper

%%%%%%%%%%%%%%%%%%%%%%%%%%%%%%%%%%%%%%%%%%%%%%%%%%%%%%%%%%%%%%%%%
% USAGE EXAMPLES
%%%%%%%%%%%%%%%%%%%%%%%%%%%%%%%%%%%%%%%%%%%%%%%%%%%%%%%%%%%%%%%%%
%
% GOOD:
%   \newcommand{\hypogenic}{\textsc{HypoGeniC}\xspace}
%   Usage: "We propose \hypogenic to generate hypotheses."
%   Result: "We propose HypoGeniC to generate hypotheses."
%
% GOOD:
%   \newcommand{\ours}{{\bf HypoGeniC}\xspace}
%   Usage: "Our method, \ours, outperforms baselines."
%   Result: "Our method, HypoGeniC, outperforms baselines."
%
% BAD (no \xspace):
%   \newcommand{\hypogenic}{\textsc{HypoGeniC}}
%   Usage: "\hypogenic outperforms..."
%   Result: "HypoGeniCoutperforms..." (missing space!)
%
%%%%%%%%%%%%%%%%%%%%%%%%%%%%%%%%%%%%%%%%%%%%%%%%%%%%%%%%%%%%%%%%%

%%%%%%%%%%%%%%%%%%%%%%%%%%%%%%%%%%%%%%%%%%%%%%%%%%%%%%%%%%%%%%%%%
% THIS PAPER
%%%%%%%%%%%%%%%%%%%%%%%%%%%%%%%%%%%%%%%%%%%%%%%%%%%%%%%%%%%%%%%%%

% Models
\newcommand{\llama}{\textsc{Llama-3.1-8B-Instruct}\xspace}
\newcommand{\llamashort}{\textsc{Llama}\xspace}
\newcommand{\qwen}{\textsc{Qwen2.5-7B-Instruct}\xspace}
\newcommand{\qwenshort}{\textsc{Qwen}\xspace}

% Edit arms
\newcommand{\icl}{\textsc{In-context}\xspace}
\newcommand{\gate}{\textsc{Gate}\xspace}
\newcommand{\rome}{\textsc{ROME-style}\xspace}
\newcommand{\ftmlp}{\textsc{FT-MLP}\xspace}
\newcommand{\loraexact}{\textsc{LoRA-exact}\xspace}
\newcommand{\loraanchor}{\textsc{LoRA+anchors}\xspace}
\newcommand{\lorakl}{\textsc{LoRA+KL}\xspace}
\newcommand{\lorapara}{\textsc{LoRA+para}\xspace}
\newcommand{\loraparakl}{\textsc{LoRA+para+KL}\xspace}
\newcommand{\loradeep}{\textsc{LoRA+para+ent+KL}\xspace}
\newcommand{\sdf}{\textsc{SDF}\xspace}

% Datasets
\newcommand{\counterfact}{\textsc{CounterFact}\xspace}
\newcommand{\triviaqa}{\textsc{TriviaQA}\xspace}
\newcommand{\wikitext}{\textsc{WikiText-2}\xspace}
\newcommand{\gsm}{\textsc{GSM8K}\xspace}

% Notation
\newcommand{\prompt}[1]{\texttt{#1}}
\newcommand{\edit}[2]{\texttt{#1}$\to$\texttt{#2}}
\newcommand{\pnew}{p(\mathrm{new})}

% Abbreviations
\newcommand{\ie}{i.e.,\xspace}
\newcommand{\eg}{e.g.,\xspace}

%
% INSTRUCTIONS:
% 1. Use \textsc{} for method/dataset/model names (small caps)
% 2. Add \xspace at the end for automatic spacing
% 3. Group related terms together
% 4. Use consistent naming conventions

\usepackage{xspace}

%%%%%%%%%%%%%%%%%%%%%%%%%%%%%%%%%%%%%%%%%%%%%%%%%%%%%%%%%%%%%%%%%
% METHOD NAMES
%%%%%%%%%%%%%%%%%%%%%%%%%%%%%%%%%%%%%%%%%%%%%%%%%%%%%%%%%%%%%%%%%

% Your method name (replace with actual name)
% Example: \newcommand{\ours}{{\bf YourMethod}\xspace}
% Example: \newcommand{\methodname}{\textsc{MethodName}\xspace}

%%%%%%%%%%%%%%%%%%%%%%%%%%%%%%%%%%%%%%%%%%%%%%%%%%%%%%%%%%%%%%%%%
% DATASET NAMES
%%%%%%%%%%%%%%%%%%%%%%%%%%%%%%%%%%%%%%%%%%%%%%%%%%%%%%%%%%%%%%%%%

% Dataset names (use \textsc for consistency)
% Example: \newcommand{\datasetone}{\textsc{Dataset-One}\xspace}
% Example: \newcommand{\datasettwo}{\textsc{Dataset-Two}\xspace}

%%%%%%%%%%%%%%%%%%%%%%%%%%%%%%%%%%%%%%%%%%%%%%%%%%%%%%%%%%%%%%%%%
% MODEL NAMES
%%%%%%%%%%%%%%%%%%%%%%%%%%%%%%%%%%%%%%%%%%%%%%%%%%%%%%%%%%%%%%%%%

% LLM and model names
% Example: \newcommand{\gpt}{\textsc{GPT-4o-mini}\xspace}
% Example: \newcommand{\llama}{\textsc{Llama-3.1-70B-Instruct}\xspace}
% Example: \newcommand{\claude}{\textsc{Claude-3.5-Sonnet}\xspace}

%%%%%%%%%%%%%%%%%%%%%%%%%%%%%%%%%%%%%%%%%%%%%%%%%%%%%%%%%%%%%%%%%
% BASELINE NAMES
%%%%%%%%%%%%%%%%%%%%%%%%%%%%%%%%%%%%%%%%%%%%%%%%%%%%%%%%%%%%%%%%%

% Baseline method names
% Example: \newcommand{\baseline}{\textsc{BaselineMethod}\xspace}
% Example: \newcommand{\fewshot}{\textsc{Few-shot}\xspace}

%%%%%%%%%%%%%%%%%%%%%%%%%%%%%%%%%%%%%%%%%%%%%%%%%%%%%%%%%%%%%%%%%
% ABBREVIATIONS
%%%%%%%%%%%%%%%%%%%%%%%%%%%%%%%%%%%%%%%%%%%%%%%%%%%%%%%%%%%%%%%%%

% Common abbreviations used in your paper
% Example: \newcommand{\ie}{i.e.,\xspace}
% Example: \newcommand{\eg}{e.g.,\xspace}
% Example: \newcommand{\cf}{cf.\xspace}
% Example: \newcommand{\wrt}{w.r.t.\xspace}

%%%%%%%%%%%%%%%%%%%%%%%%%%%%%%%%%%%%%%%%%%%%%%%%%%%%%%%%%%%%%%%%%
% DOMAIN-SPECIFIC TERMS
%%%%%%%%%%%%%%%%%%%%%%%%%%%%%%%%%%%%%%%%%%%%%%%%%%%%%%%%%%%%%%%%%

% Terms specific to your research domain
% Keep consistent formatting throughout the paper

%%%%%%%%%%%%%%%%%%%%%%%%%%%%%%%%%%%%%%%%%%%%%%%%%%%%%%%%%%%%%%%%%
% USAGE EXAMPLES
%%%%%%%%%%%%%%%%%%%%%%%%%%%%%%%%%%%%%%%%%%%%%%%%%%%%%%%%%%%%%%%%%
%
% GOOD:
%   \newcommand{\hypogenic}{\textsc{HypoGeniC}\xspace}
%   Usage: "We propose \hypogenic to generate hypotheses."
%   Result: "We propose HypoGeniC to generate hypotheses."
%
% GOOD:
%   \newcommand{\ours}{{\bf HypoGeniC}\xspace}
%   Usage: "Our method, \ours, outperforms baselines."
%   Result: "Our method, HypoGeniC, outperforms baselines."
%
% BAD (no \xspace):
%   \newcommand{\hypogenic}{\textsc{HypoGeniC}}
%   Usage: "\hypogenic outperforms..."
%   Result: "HypoGeniCoutperforms..." (missing space!)
%
%%%%%%%%%%%%%%%%%%%%%%%%%%%%%%%%%%%%%%%%%%%%%%%%%%%%%%%%%%%%%%%%%

%%%%%%%%%%%%%%%%%%%%%%%%%%%%%%%%%%%%%%%%%%%%%%%%%%%%%%%%%%%%%%%%%
% THIS PAPER
%%%%%%%%%%%%%%%%%%%%%%%%%%%%%%%%%%%%%%%%%%%%%%%%%%%%%%%%%%%%%%%%%

% Models
\newcommand{\llama}{\textsc{Llama-3.1-8B-Instruct}\xspace}
\newcommand{\llamashort}{\textsc{Llama}\xspace}
\newcommand{\qwen}{\textsc{Qwen2.5-7B-Instruct}\xspace}
\newcommand{\qwenshort}{\textsc{Qwen}\xspace}

% Edit arms
\newcommand{\icl}{\textsc{In-context}\xspace}
\newcommand{\gate}{\textsc{Gate}\xspace}
\newcommand{\rome}{\textsc{ROME-style}\xspace}
\newcommand{\ftmlp}{\textsc{FT-MLP}\xspace}
\newcommand{\loraexact}{\textsc{LoRA-exact}\xspace}
\newcommand{\loraanchor}{\textsc{LoRA+anchors}\xspace}
\newcommand{\lorakl}{\textsc{LoRA+KL}\xspace}
\newcommand{\lorapara}{\textsc{LoRA+para}\xspace}
\newcommand{\loraparakl}{\textsc{LoRA+para+KL}\xspace}
\newcommand{\loradeep}{\textsc{LoRA+para+ent+KL}\xspace}
\newcommand{\sdf}{\textsc{SDF}\xspace}

% Datasets
\newcommand{\counterfact}{\textsc{CounterFact}\xspace}
\newcommand{\triviaqa}{\textsc{TriviaQA}\xspace}
\newcommand{\wikitext}{\textsc{WikiText-2}\xspace}
\newcommand{\gsm}{\textsc{GSM8K}\xspace}

% Notation
\newcommand{\prompt}[1]{\texttt{#1}}
\newcommand{\edit}[2]{\texttt{#1}$\to$\texttt{#2}}
\newcommand{\pnew}{p(\mathrm{new})}

% Abbreviations
\newcommand{\ie}{i.e.,\xspace}
\newcommand{\eg}{e.g.,\xspace}

%
% INSTRUCTIONS:
% 1. Use \textsc{} for method/dataset/model names (small caps)
% 2. Add \xspace at the end for automatic spacing
% 3. Group related terms together
% 4. Use consistent naming conventions

\usepackage{xspace}

%%%%%%%%%%%%%%%%%%%%%%%%%%%%%%%%%%%%%%%%%%%%%%%%%%%%%%%%%%%%%%%%%
% METHOD NAMES
%%%%%%%%%%%%%%%%%%%%%%%%%%%%%%%%%%%%%%%%%%%%%%%%%%%%%%%%%%%%%%%%%

% Your method name (replace with actual name)
% Example: \newcommand{\ours}{{\bf YourMethod}\xspace}
% Example: \newcommand{\methodname}{\textsc{MethodName}\xspace}

%%%%%%%%%%%%%%%%%%%%%%%%%%%%%%%%%%%%%%%%%%%%%%%%%%%%%%%%%%%%%%%%%
% DATASET NAMES
%%%%%%%%%%%%%%%%%%%%%%%%%%%%%%%%%%%%%%%%%%%%%%%%%%%%%%%%%%%%%%%%%

% Dataset names (use \textsc for consistency)
% Example: \newcommand{\datasetone}{\textsc{Dataset-One}\xspace}
% Example: \newcommand{\datasettwo}{\textsc{Dataset-Two}\xspace}

%%%%%%%%%%%%%%%%%%%%%%%%%%%%%%%%%%%%%%%%%%%%%%%%%%%%%%%%%%%%%%%%%
% MODEL NAMES
%%%%%%%%%%%%%%%%%%%%%%%%%%%%%%%%%%%%%%%%%%%%%%%%%%%%%%%%%%%%%%%%%

% LLM and model names
% Example: \newcommand{\gpt}{\textsc{GPT-4o-mini}\xspace}
% Example: \newcommand{\llama}{\textsc{Llama-3.1-70B-Instruct}\xspace}
% Example: \newcommand{\claude}{\textsc{Claude-3.5-Sonnet}\xspace}

%%%%%%%%%%%%%%%%%%%%%%%%%%%%%%%%%%%%%%%%%%%%%%%%%%%%%%%%%%%%%%%%%
% BASELINE NAMES
%%%%%%%%%%%%%%%%%%%%%%%%%%%%%%%%%%%%%%%%%%%%%%%%%%%%%%%%%%%%%%%%%

% Baseline method names
% Example: \newcommand{\baseline}{\textsc{BaselineMethod}\xspace}
% Example: \newcommand{\fewshot}{\textsc{Few-shot}\xspace}

%%%%%%%%%%%%%%%%%%%%%%%%%%%%%%%%%%%%%%%%%%%%%%%%%%%%%%%%%%%%%%%%%
% ABBREVIATIONS
%%%%%%%%%%%%%%%%%%%%%%%%%%%%%%%%%%%%%%%%%%%%%%%%%%%%%%%%%%%%%%%%%

% Common abbreviations used in your paper
% Example: \newcommand{\ie}{i.e.,\xspace}
% Example: \newcommand{\eg}{e.g.,\xspace}
% Example: \newcommand{\cf}{cf.\xspace}
% Example: \newcommand{\wrt}{w.r.t.\xspace}

%%%%%%%%%%%%%%%%%%%%%%%%%%%%%%%%%%%%%%%%%%%%%%%%%%%%%%%%%%%%%%%%%
% DOMAIN-SPECIFIC TERMS
%%%%%%%%%%%%%%%%%%%%%%%%%%%%%%%%%%%%%%%%%%%%%%%%%%%%%%%%%%%%%%%%%

% Terms specific to your research domain
% Keep consistent formatting throughout the paper

%%%%%%%%%%%%%%%%%%%%%%%%%%%%%%%%%%%%%%%%%%%%%%%%%%%%%%%%%%%%%%%%%
% USAGE EXAMPLES
%%%%%%%%%%%%%%%%%%%%%%%%%%%%%%%%%%%%%%%%%%%%%%%%%%%%%%%%%%%%%%%%%
%
% GOOD:
%   \newcommand{\hypogenic}{\textsc{HypoGeniC}\xspace}
%   Usage: "We propose \hypogenic to generate hypotheses."
%   Result: "We propose HypoGeniC to generate hypotheses."
%
% GOOD:
%   \newcommand{\ours}{{\bf HypoGeniC}\xspace}
%   Usage: "Our method, \ours, outperforms baselines."
%   Result: "Our method, HypoGeniC, outperforms baselines."
%
% BAD (no \xspace):
%   \newcommand{\hypogenic}{\textsc{HypoGeniC}}
%   Usage: "\hypogenic outperforms..."
%   Result: "HypoGeniCoutperforms..." (missing space!)
%
%%%%%%%%%%%%%%%%%%%%%%%%%%%%%%%%%%%%%%%%%%%%%%%%%%%%%%%%%%%%%%%%%

%%%%%%%%%%%%%%%%%%%%%%%%%%%%%%%%%%%%%%%%%%%%%%%%%%%%%%%%%%%%%%%%%
% THIS PAPER
%%%%%%%%%%%%%%%%%%%%%%%%%%%%%%%%%%%%%%%%%%%%%%%%%%%%%%%%%%%%%%%%%

% Models
\newcommand{\llama}{\textsc{Llama-3.1-8B-Instruct}\xspace}
\newcommand{\llamashort}{\textsc{Llama}\xspace}
\newcommand{\qwen}{\textsc{Qwen2.5-7B-Instruct}\xspace}
\newcommand{\qwenshort}{\textsc{Qwen}\xspace}

% Edit arms
\newcommand{\icl}{\textsc{In-context}\xspace}
\newcommand{\gate}{\textsc{Gate}\xspace}
\newcommand{\rome}{\textsc{ROME-style}\xspace}
\newcommand{\ftmlp}{\textsc{FT-MLP}\xspace}
\newcommand{\loraexact}{\textsc{LoRA-exact}\xspace}
\newcommand{\loraanchor}{\textsc{LoRA+anchors}\xspace}
\newcommand{\lorakl}{\textsc{LoRA+KL}\xspace}
\newcommand{\lorapara}{\textsc{LoRA+para}\xspace}
\newcommand{\loraparakl}{\textsc{LoRA+para+KL}\xspace}
\newcommand{\loradeep}{\textsc{LoRA+para+ent+KL}\xspace}
\newcommand{\sdf}{\textsc{SDF}\xspace}

% Datasets
\newcommand{\counterfact}{\textsc{CounterFact}\xspace}
\newcommand{\triviaqa}{\textsc{TriviaQA}\xspace}
\newcommand{\wikitext}{\textsc{WikiText-2}\xspace}
\newcommand{\gsm}{\textsc{GSM8K}\xspace}

% Notation
\newcommand{\prompt}[1]{\texttt{#1}}
\newcommand{\edit}[2]{\texttt{#1}$\to$\texttt{#2}}
\newcommand{\pnew}{p(\mathrm{new})}

% Abbreviations
\newcommand{\ie}{i.e.,\xspace}
\newcommand{\eg}{e.g.,\xspace}

%
% INSTRUCTIONS:
% 1. Use \textsc{} for method/dataset/model names (small caps)
% 2. Add \xspace at the end for automatic spacing
% 3. Group related terms together
% 4. Use consistent naming conventions

\usepackage{xspace}

%%%%%%%%%%%%%%%%%%%%%%%%%%%%%%%%%%%%%%%%%%%%%%%%%%%%%%%%%%%%%%%%%
% METHOD NAMES
%%%%%%%%%%%%%%%%%%%%%%%%%%%%%%%%%%%%%%%%%%%%%%%%%%%%%%%%%%%%%%%%%

% Your method name (replace with actual name)
% Example: \newcommand{\ours}{{\bf YourMethod}\xspace}
% Example: \newcommand{\methodname}{\textsc{MethodName}\xspace}

%%%%%%%%%%%%%%%%%%%%%%%%%%%%%%%%%%%%%%%%%%%%%%%%%%%%%%%%%%%%%%%%%
% DATASET NAMES
%%%%%%%%%%%%%%%%%%%%%%%%%%%%%%%%%%%%%%%%%%%%%%%%%%%%%%%%%%%%%%%%%

% Dataset names (use \textsc for consistency)
% Example: \newcommand{\datasetone}{\textsc{Dataset-One}\xspace}
% Example: \newcommand{\datasettwo}{\textsc{Dataset-Two}\xspace}

%%%%%%%%%%%%%%%%%%%%%%%%%%%%%%%%%%%%%%%%%%%%%%%%%%%%%%%%%%%%%%%%%
% MODEL NAMES
%%%%%%%%%%%%%%%%%%%%%%%%%%%%%%%%%%%%%%%%%%%%%%%%%%%%%%%%%%%%%%%%%

% LLM and model names
% Example: \newcommand{\gpt}{\textsc{GPT-4o-mini}\xspace}
% Example: \newcommand{\llama}{\textsc{Llama-3.1-70B-Instruct}\xspace}
% Example: \newcommand{\claude}{\textsc{Claude-3.5-Sonnet}\xspace}

%%%%%%%%%%%%%%%%%%%%%%%%%%%%%%%%%%%%%%%%%%%%%%%%%%%%%%%%%%%%%%%%%
% BASELINE NAMES
%%%%%%%%%%%%%%%%%%%%%%%%%%%%%%%%%%%%%%%%%%%%%%%%%%%%%%%%%%%%%%%%%

% Baseline method names
% Example: \newcommand{\baseline}{\textsc{BaselineMethod}\xspace}
% Example: \newcommand{\fewshot}{\textsc{Few-shot}\xspace}

%%%%%%%%%%%%%%%%%%%%%%%%%%%%%%%%%%%%%%%%%%%%%%%%%%%%%%%%%%%%%%%%%
% ABBREVIATIONS
%%%%%%%%%%%%%%%%%%%%%%%%%%%%%%%%%%%%%%%%%%%%%%%%%%%%%%%%%%%%%%%%%

% Common abbreviations used in your paper
% Example: \newcommand{\ie}{i.e.,\xspace}
% Example: \newcommand{\eg}{e.g.,\xspace}
% Example: \newcommand{\cf}{cf.\xspace}
% Example: \newcommand{\wrt}{w.r.t.\xspace}

%%%%%%%%%%%%%%%%%%%%%%%%%%%%%%%%%%%%%%%%%%%%%%%%%%%%%%%%%%%%%%%%%
% DOMAIN-SPECIFIC TERMS
%%%%%%%%%%%%%%%%%%%%%%%%%%%%%%%%%%%%%%%%%%%%%%%%%%%%%%%%%%%%%%%%%

% Terms specific to your research domain
% Keep consistent formatting throughout the paper

%%%%%%%%%%%%%%%%%%%%%%%%%%%%%%%%%%%%%%%%%%%%%%%%%%%%%%%%%%%%%%%%%
% USAGE EXAMPLES
%%%%%%%%%%%%%%%%%%%%%%%%%%%%%%%%%%%%%%%%%%%%%%%%%%%%%%%%%%%%%%%%%
%
% GOOD:
%   \newcommand{\hypogenic}{\textsc{HypoGeniC}\xspace}
%   Usage: "We propose \hypogenic to generate hypotheses."
%   Result: "We propose HypoGeniC to generate hypotheses."
%
% GOOD:
%   \newcommand{\ours}{{\bf HypoGeniC}\xspace}
%   Usage: "Our method, \ours, outperforms baselines."
%   Result: "Our method, HypoGeniC, outperforms baselines."
%
% BAD (no \xspace):
%   \newcommand{\hypogenic}{\textsc{HypoGeniC}}
%   Usage: "\hypogenic outperforms..."
%   Result: "HypoGeniCoutperforms..." (missing space!)
%
%%%%%%%%%%%%%%%%%%%%%%%%%%%%%%%%%%%%%%%%%%%%%%%%%%%%%%%%%%%%%%%%%

%%%%%%%%%%%%%%%%%%%%%%%%%%%%%%%%%%%%%%%%%%%%%%%%%%%%%%%%%%%%%%%%%
% THIS PAPER
%%%%%%%%%%%%%%%%%%%%%%%%%%%%%%%%%%%%%%%%%%%%%%%%%%%%%%%%%%%%%%%%%

% Models
\newcommand{\llama}{\textsc{Llama-3.1-8B-Instruct}\xspace}
\newcommand{\llamashort}{\textsc{Llama}\xspace}
\newcommand{\qwen}{\textsc{Qwen2.5-7B-Instruct}\xspace}
\newcommand{\qwenshort}{\textsc{Qwen}\xspace}

% Edit arms
\newcommand{\icl}{\textsc{In-context}\xspace}
\newcommand{\gate}{\textsc{Gate}\xspace}
\newcommand{\rome}{\textsc{ROME-style}\xspace}
\newcommand{\ftmlp}{\textsc{FT-MLP}\xspace}
\newcommand{\loraexact}{\textsc{LoRA-exact}\xspace}
\newcommand{\loraanchor}{\textsc{LoRA+anchors}\xspace}
\newcommand{\lorakl}{\textsc{LoRA+KL}\xspace}
\newcommand{\lorapara}{\textsc{LoRA+para}\xspace}
\newcommand{\loraparakl}{\textsc{LoRA+para+KL}\xspace}
\newcommand{\loradeep}{\textsc{LoRA+para+ent+KL}\xspace}
\newcommand{\sdf}{\textsc{SDF}\xspace}

% Datasets
\newcommand{\counterfact}{\textsc{CounterFact}\xspace}
\newcommand{\triviaqa}{\textsc{TriviaQA}\xspace}
\newcommand{\wikitext}{\textsc{WikiText-2}\xspace}
\newcommand{\gsm}{\textsc{GSM8K}\xspace}

% Notation
\newcommand{\prompt}[1]{\texttt{#1}}
\newcommand{\edit}[2]{\texttt{#1}$\to$\texttt{#2}}
\newcommand{\pnew}{p(\mathrm{new})}

% Abbreviations
\newcommand{\ie}{i.e.,\xspace}
\newcommand{\eg}{e.g.,\xspace}
